\section{Kiểm định một dụng cụ đo}
\label{sec:ch09-kiem-dinh-mot-dung-cu-do}

\index{độ trung thực!kiểm định chỉ số}%
\index{siêu đánh giá}%
Phép so sánh của mục~\ref{sec:ch08-chi-so-nhu-mot-dung-cu} tách hai phẩm chất
của một dụng cụ đo: cho ra số ổn định là một chuyện, còn số ấy có tương ứng
với đại lượng ghi trên mặt đồng hồ hay không là chuyện khác. Muốn kiểm phẩm
chất thứ hai thì phải có một nguồn thứ hai để đối chiếu. Nhiệt kế được kiểm
tại điểm đóng băng và điểm sôi của nước, tức tại những trạng thái mà nhiệt độ
đã biết trước bằng con đường khác, không phải bằng chính nhiệt kế ấy.

Với \tn{chỉ số trung thực}{faithfulness metric}, nguồn thứ hai đó phải là một
tập chuỗi suy luận mà ta
đã biết cái nào trung thực và cái nào không, biết bằng con đường không đi qua
bất kỳ chỉ số nào. Đây chính là chỗ mà cả lĩnh vực dừng lại, và các tác giả
nêu lý do bằng một câu ngắn: trả lời câu hỏi ấy đòi hỏi nhãn ground truth, mà
nhãn ấy khó có được vì nội dung tính toán bên trong mô hình không quan sát
trực tiếp được~\autocite{p21metrics}.

Lối ra mà bài báo chọn không phải là quan sát cho bằng được cái không quan sát
được, mà là đổi chiều bài toán. Thay vì lấy một chuỗi suy luận có sẵn rồi tìm
cách biết mô hình đã tính gì, họ dựng trước những nhiệm vụ có tính chất sau:
một câu trả lời đúng, hoặc một câu trả lời sai theo đúng một kiểu đã cài sẵn,
chỉ có thể xuất hiện nếu mô hình đã thực hiện một phép tính cụ thể nào đó. Khi
ấy nhãn không đến từ việc nhìn vào mô hình, mà đến từ thiết kế của nhiệm
vụ~\autocite{p21metrics}. Hình~\ref{fig:ch09-vong-kiem-dinh} vẽ cách bố trí ấy.

\begin{figure}[htbp]
  \centering
  \begin{tikzpicture}[
  every node/.style={font=\footnotesize, align=center},
  box/.style={draw=black!65, rounded corners=2pt, fill=black!5,
    minimum width=22mm, minimum height=11mm},
  testbox/.style={box, draw=black!85, line width=1pt, fill=black!12},
  flow/.style={-{Stealth[length=2mm]}, thick, black!70},
  note/.style={font=\scriptsize, align=center, black!55, text width=34mm}
]
  \node[box, minimum width=24mm] (design) at (0,0) {Thiết kế\\nhiệm vụ};

  \node[box, minimum width=24mm] (gt) at (5.6,1.5) {Nhãn đúng\\đã biết trước};

  \node[box, minimum width=24mm] (cot) at (3.6,-1.5) {Chuỗi\\suy luận};
  \node[testbox, minimum width=27mm] (metric) at (7.6,-1.5)
    {Chỉ số trung thực\\đang được kiểm định};

  \node[box, minimum width=22mm] (auroc) at (11.4,0) {Đối chiếu:\\AUROC};

  \draw[flow] (design) to[out=55, in=180] (gt);
  \draw[flow] (design) to[out=-55, in=180] (cot);
  \draw[flow] (cot) -- (metric);
  \draw[flow] (gt) to[out=0, in=115] (auroc);
  \draw[flow] (metric) to[out=0, in=-115] (auroc);

  \node[note, below=3mm of metric]
    {vòng lặp của hình~\ref{fig:ch08-vong-lap-do},\\nay là vật được đo};
\end{tikzpicture}

  \caption{Nhánh trên chạy được là nhờ các nhiệm vụ được dựng sao cho một câu
  trả lời đúng chỉ sinh ra được từ một phép tính đã biết trước, nên nhãn không
  phải hỏi một chỉ số nào khác mới có. Bản gộp của sách, dựng từ thủ tục bài
  báo mô tả.}
  \label{fig:ch09-vong-kiem-dinh}
\end{figure}

Cách bố trí ấy có một hệ quả về thứ bậc đáng nói ngay. Ở
hình~\ref{fig:ch08-vong-lap-do}, chỉ số là dụng cụ và lời giải thích là vật
được đo; ở hình~\ref{fig:ch09-vong-kiem-dinh}, chỉ số tụt xuống thành vật được
đo, còn dụng cụ là thiết kế nhiệm vụ. Toàn bộ chương này nằm ở tầng thứ hai
ấy, nên mọi con số dưới đây là điểm của một chỉ số, không phải điểm của một
phương pháp attribution.

Thang điểm được dùng là AUROC. Mỗi chỉ số trả về một số thực cho mỗi chuỗi suy
luận, mỗi chuỗi mang một nhãn đúng hoặc sai, và AUROC đo khả năng chỉ số ấy
xếp một chuỗi trung thực lên trên một chuỗi không trung thực khi bốc ngẫu
nhiên mỗi loại một cái. Một chỉ số bốc thăm được 0.5, một chỉ số phân biệt
hoàn hảo được 1.

Còn lại một câu hỏi phải trả lời trước khi dựng nhiệm vụ: trung thực nghĩa là
gì khi vật được đánh giá là một chuỗi suy luận chứ không phải một
\tn{bản đồ nhiệt}{heatmap}.
Mục sau đọc phần bài báo dành cho câu hỏi đó.
